\documentclass[11pt]{article}

\usepackage[margin=1in]{geometry}
\usepackage{amsmath,amssymb}
\usepackage{array}
\usepackage{enumitem}

\title{Lecture 4 Notes: \\ Near Opposites: Kant and Mill}
\author{Philosophy of Mathematics}
\date{}

\begin{document}

\maketitle

\section*{Central Theme}

Kant and Mill are near opposites because both reject simple Platonism, but they explain
mathematics in almost opposite ways.

\[
\boxed{
\text{Kant: mathematics is synthetic a priori.}
}
\]

\[
\boxed{
\text{Mill: mathematics is synthetic a posteriori.}
}
\]

The guiding contrast:

\[
\begin{array}{c|c}
\textbf{Kant} & \textbf{Mill} \\
\hline
\text{mathematics grounded in pure intuition} &
\text{mathematics grounded in experience} \\
\text{necessary and a priori} &
\text{empirical and inductive} \\
\text{forms of possible experience} &
\text{generalizations from actual experience}
\end{array}
\]

\section*{1. Reorientation: Mathematics after the Scientific Revolution}

By Kant's time, mathematics and science had changed dramatically.

The seventeenth century brought:

\[
\text{analytic geometry}
\]

\[
\text{calculus}
\]

\[
\text{Newtonian mechanics}
\]

\[
\text{new ideas of continuity, derivative, and limit}
\]

Mathematics was no longer only geometry and arithmetic in the ancient sense. It had
become deeply connected to modern physics.

\subsection*{Teaching Point}

The philosophy of mathematics now has to explain not only ancient geometry and
arithmetic, but also the role of mathematics in modern science.

\subsection*{Whiteboard}

\[
\text{Descartes}
\quad
\text{Newton}
\quad
\text{Leibniz}
\]

\[
\text{geometry} + \text{algebra} + \text{calculus} + \text{physics}
\]

\[
\text{New question:}
\quad
\text{How does mathematics make modern science possible?}
\]

\section*{2. Rationalism and Empiricism}

The background dispute is between rationalism and empiricism.

\[
\begin{array}{c|c}
\textbf{Rationalism} & \textbf{Empiricism} \\
\hline
\text{reason as source of knowledge} &
\text{sense experience as source of knowledge} \\
\text{mathematics as model} &
\text{mathematics as problem} \\
\text{Descartes, Spinoza, Leibniz} &
\text{Locke, Berkeley, Hume, Reid}
\end{array}
\]

Rationalists explain the necessity of mathematics well. Mathematical truths are grasped
by reason and do not depend on contingent experience.

Empiricists explain the applicability of mathematics well. Mathematical ideas arise from
experience with physical objects.

But each view has a problem.

\[
\begin{array}{c|c}
\textbf{Rationalism} & \textbf{Empiricism} \\
\hline
\text{necessity explained} & \text{applicability explained} \\
\text{applicability mysterious} & \text{necessity mysterious}
\end{array}
\]

\subsection*{Teaching Point}

Kant and Mill inherit the old Plato/Aristotle tension: exactness and necessity pull us
toward rationalism, while applicability pulls us toward empiricism.

\subsection*{Whiteboard}

\[
\text{Problem for rationalism:}
\quad
\text{Why does pure reason apply to nature?}
\]

\[
\text{Problem for empiricism:}
\quad
\text{Why does mathematics seem necessary?}
\]

\section*{3. Kant's Central Problem}

Kant tries to synthesize rationalism and empiricism.

His problem is:

\[
\boxed{
\text{How can mathematics be a priori and yet apply universally to experience?}
}
\]

He wants to explain three features at once:

\[
\text{necessity}
\quad + \quad
\text{a priority}
\quad + \quad
\text{applicability}
\]

Kant's answer is that mathematics is \emph{synthetic a priori} and grounded in
\emph{intuition}.

\subsection*{Teaching Point}

Kant's philosophy of mathematics is an attempt to preserve the necessity of mathematics
without making mathematics irrelevant to empirical science.

\subsection*{Whiteboard}

\[
\textbf{Kant's task:}
\]

\[
\text{necessary}
+
\text{a priori}
+
\text{applicable}
\]

\[
\boxed{
\text{Mathematics} = \text{synthetic a priori}
}
\]

\section*{4. Analytic and Synthetic}

Kant distinguishes analytic from synthetic judgements.

\[
\begin{array}{c|c}
\textbf{Analytic} & \textbf{Synthetic} \\
\hline
\text{predicate contained in subject concept} &
\text{predicate adds something new} \\
\text{known by conceptual analysis} &
\text{requires intuition or experience} \\
\text{elucidatory} &
\text{ampliative}
\end{array}
\]

Example of an analytic judgement:

\[
\text{All bachelors are unmarried.}
\]

The predicate concept ``unmarried'' is already contained in the subject concept
``bachelor.''

Kant thinks most mathematical judgements are not like this. They are synthetic.

\subsection*{Teaching Point}

For Kant, mathematical propositions do not merely unpack definitions. They extend our
knowledge.

\subsection*{Whiteboard}

\[
\text{Analytic:}
\quad
\text{unpack the concept}
\]

\[
\text{Synthetic:}
\quad
\text{add something new}
\]

\[
\text{Question:}
\quad
7+5=12
\quad
\text{analytic or synthetic?}
\]

\section*{5. Kant on Arithmetic: \(7+5=12\)}

Kant famously argues that:

\[
7+5=12
\]

is synthetic.

The concept of the sum of seven and five contains only the idea of combining seven and
five into one number. It does not already contain the concept of twelve.

To get twelve, we must construct the sum.

\[
7
\quad \rightarrow \quad
8
\quad \rightarrow \quad
9
\quad \rightarrow \quad
10
\quad \rightarrow \quad
11
\quad \rightarrow \quad
12
\]

This construction involves intuition, especially succession in time.

\subsection*{Teaching Point}

For Kant, arithmetic is not mere definition. It requires constructive activity in intuition.

\subsection*{Whiteboard}

\[
7+5
\neq
\text{conceptual analysis alone}
\]

\[
7,8,9,10,11,12
\]

\[
\text{Arithmetic uses succession.}
\]

\section*{6. Kant on Geometry: Construction in Intuition}

Kant also thinks geometry is synthetic a priori.

A philosopher who only analyzes the concept of triangle will not discover that the angles
of a triangle sum to two right angles.

The geometer does something different. She constructs a triangle and then adds auxiliary
constructions.

For example, in Euclid's proof that the angles of a triangle sum to two right angles, the
geometer:

\[
\text{constructs a triangle}
\]

\[
\text{extends one side}
\]

\[
\text{draws a parallel line}
\]

\[
\text{derives the angle sum}
\]

\subsection*{Teaching Point}

Kant's model of mathematics is not passive inspection of concepts. It is active
construction in pure intuition.

\subsection*{Whiteboard}

\[
\text{Concept of triangle}
\quad
\not\Rightarrow
\quad
\text{angle sum } = 180^\circ
\]

\[
\text{construct}
\quad \rightarrow \quad
\text{extend}
\quad \rightarrow \quad
\text{draw parallel}
\quad \rightarrow \quad
\text{prove}
\]

\section*{7. Philosophy and Mathematics}

Kant contrasts philosophical cognition with mathematical cognition.

\[
\begin{array}{c|c}
\textbf{Philosophy} & \textbf{Mathematics} \\
\hline
\text{rational cognition from concepts} &
\text{rational cognition from construction of concepts} \\
\text{conceptual analysis} &
\text{intuition and construction} \\
\text{makes explicit what is implicit} &
\text{produces new knowledge}
\end{array}
\]

This is why mathematics can be synthetic and yet a priori.

It is not based on empirical observation, but it is also not based merely on definitions.

\subsection*{Teaching Point}

Kant gives mathematics its own distinctive epistemology: neither ordinary empirical
knowledge nor mere conceptual analysis.

\subsection*{Whiteboard}

\[
\text{Philosophy: concepts}
\]

\[
\text{Mathematics: construction of concepts}
\]

\[
\text{Mathematics is not just logic or definition.}
\]

\section*{8. Pure Intuition: Space and Time}

Kant's key idea is that mathematics is grounded in pure intuition.

\[
\text{Geometry}
\quad \rightarrow \quad
\text{space}
\]

\[
\text{Arithmetic}
\quad \rightarrow \quad
\text{time/succession}
\]

Space and time are not learned from ordinary experience. Rather, they are forms through
which we experience objects.

Mathematics applies to experience because experience itself must be structured spatially
and temporally.

\subsection*{Teaching Point}

Kant's great move is to locate mathematics neither in a Platonic realm nor in ordinary
empirical objects, but in the form of possible experience.

\subsection*{Whiteboard}

\[
\text{Mathematics concerns the form of experience.}
\]

\[
\text{Geometry} = \text{form of space}
\]

\[
\text{Arithmetic} = \text{form of succession}
\]

\[
\text{That is why mathematics applies.}
\]

\section*{9. Kant's Strength}

Kant's view explains why mathematics seems necessary, a priori, and applicable.

\[
\begin{array}{c|c}
\textbf{Feature} & \textbf{Kant's Explanation} \\
\hline
\text{necessary} &
\text{experience must conform to space and time} \\
\text{a priori} &
\text{space and time are forms of intuition, not empirical discoveries} \\
\text{applicable} &
\text{objects of experience appear within those forms}
\end{array}
\]

Kant therefore tries to solve the old problem that divided rationalists and empiricists.

\subsection*{Teaching Point}

Kant's view is powerful because it explains both the certainty of mathematics and its
application to experience.

\subsection*{Whiteboard}

\[
\textbf{Kant's success:}
\]

\[
\text{mathematics structures possible experience}
\]

\[
\therefore
\quad
\text{a priori and applicable}
\]

\section*{10. Kant's Problem: Non-Euclidean Geometry}

Kant's view is placed under pressure by later developments in mathematics and physics.

The central example is non-Euclidean geometry.

Kant held that Euclidean geometry describes the necessary form of spatial intuition. But
modern physics uses non-Euclidean geometry to describe space-time.

This creates a serious problem:

\[
\boxed{
\text{Can a supposedly a priori truth be overturned by empirical science?}
}
\]

If Euclidean geometry is synthetic a priori, then it should not be empirically revisable. But
the development of non-Euclidean geometry and its application to physics suggests
otherwise.

\subsection*{Teaching Point}

Non-Euclidean geometry threatens Kant's claim that Euclidean geometry gives the
necessary form of spatial experience.

\subsection*{Whiteboard}

\[
\text{Kant: Euclidean geometry is necessary.}
\]

\[
\text{Modern physics: space-time may be non-Euclidean.}
\]

\[
\text{Problem:}
\quad
\text{What happens to the synthetic a priori?}
\]

\section*{11. Further Problems for Kant}

Other developments in mathematics also strain Kant's account.

Examples include:

\[
\text{complex analysis}
\]

\[
\text{higher-dimensional geometry}
\]

\[
\text{functional analysis}
\]

\[
\text{set theory}
\]

These areas do not obviously concern the forms of ordinary spatial or temporal
perception.

Kant's view works best for Euclidean geometry and elementary arithmetic. It is less clear
how it extends to modern mathematics.

\subsection*{Teaching Point}

Kant gives a powerful account of traditional mathematics, but later mathematics stretches
or breaks the link between mathematics and intuition.

\subsection*{Whiteboard}

\[
\text{Euclidean geometry}
\quad \text{fits Kant well}
\]

\[
\text{set theory, complex analysis, higher dimensions}
\quad
\text{fit less well}
\]

\section*{12. Mill's Radical Empiricism}

Mill takes the opposite route from Kant.

He rejects a priori mathematical knowledge. Mathematics is not grounded in pure
intuition. It is grounded in experience.

\[
\boxed{
\text{Mathematics is empirical.}
}
\]

Mill distinguishes verbal propositions from real propositions.

\[
\begin{array}{c|c}
\textbf{Verbal propositions} & \textbf{Real propositions} \\
\hline
\text{true by definition} &
\text{about the world} \\
\text{no genuine factual content} &
\text{empirical content} \\
\text{not mathematics} &
\text{mathematics}
\end{array}
\]

For Mill, mathematics and much of logic consist of real propositions. They are ultimately
known through induction from experience.

\subsection*{Teaching Point}

Mill agrees with Kant that mathematics is not merely analytic. But he denies Kant's
a priori intuition and makes mathematics empirical.

\subsection*{Whiteboard}

\[
\textbf{Kant:}
\quad
\text{synthetic a priori}
\]

\[
\textbf{Mill:}
\quad
\text{synthetic a posteriori}
\]

\[
\text{Mathematics} = \text{real propositions about the world}
\]

\section*{13. Mill's Epistemology: Induction}

Mill's basic model of knowledge is enumerative induction.

\[
\text{many observed cases}
\quad \longrightarrow \quad
\text{general law}
\]

Example:

\[
\text{many black crows}
\quad \longrightarrow \quad
\text{all crows are black}
\]

Mill applies this model to mathematics.

Mathematical propositions summarize regularities discovered in experience. We repeatedly
combine and separate objects and observe stable patterns.

\subsection*{Teaching Point}

Mill's view is naturalistic: mathematics is part of the same empirical inquiry as science.

\subsection*{Whiteboard}

\[
\text{Observation}
\quad \longrightarrow \quad
\text{generalization}
\quad \longrightarrow \quad
\text{mathematics}
\]

\[
\text{No special a priori faculty.}
\]

\section*{14. Mill on Geometry}

Mill rejects abstract geometrical objects.

A geometrical line is not a Platonic object or a Kantian construction in pure intuition. It is
an idealization of physical lines.

\[
\text{drawn line}
\quad \rightarrow \quad
\text{thinner}
\quad \rightarrow \quad
\text{straighter}
\quad \rightarrow \quad
\text{ideal line}
\]

Points, lines, and circles are limit concepts.

Strictly speaking, Euclidean geometry is fictional, because there are no perfect
breadthless lines or dimensionless points in nature. But it is useful because real physical
objects approximate geometrical figures.

\subsection*{Teaching Point}

Mill explains the applicability of geometry by making geometry an idealized empirical
science.

\subsection*{Whiteboard}

\[
\textbf{Millian geometry:}
\]

\[
\text{idealized physical figures}
\]

\[
\text{perfect line} = \text{limit of thinner and straighter drawn lines}
\]

\[
\text{geometry is approximately true of nature}
\]

\section*{15. Mill's Problem: Exact Geometry}

Mill's geometry faces a difficulty.

If geometry is based on physical experience, how can it yield exact truths?

For example:

\[
\text{Between any two points, a straight line can be drawn.}
\]

But physically, we cannot draw a breadthless line. Nor can we bisect a line segment
indefinitely. Eventually we reach physical limits.

So Mill must appeal to idealization or possibility. But then the question becomes:

\[
\boxed{
\text{Where does this notion of possibility come from?}
}
\]

If all knowledge comes from experience, how do we know what is ideally possible beyond
experience?

\subsection*{Teaching Point}

Mill explains applicability, but at the cost of making exactness difficult to understand.

\subsection*{Whiteboard}

\[
\text{Experience gives approximate figures.}
\]

\[
\text{Geometry gives exact truths.}
\]

\[
\text{Problem:}
\quad
\text{How does approximation yield exactness?}
\]

\section*{16. Mill on Arithmetic}

Mill also treats arithmetic empirically.

Numbers are always numbers of something.

\[
\text{ten bodies}
\qquad
\text{ten sounds}
\qquad
\text{ten heartbeats}
\]

There are no numbers in the abstract.

Numerals are not names of abstract objects. Rather, they are general terms applying to
aggregates.

\[
\text{``two'' denotes all pairs}
\]

\[
\text{``twelve'' denotes all dozens}
\]

Arithmetic sums are generalizations about what happens when aggregates are collected or
separated.

\[
2+1=3
\]

means, roughly:

\[
\text{two things combined with one thing make three things.}
\]

\subsection*{Teaching Point}

Mill makes arithmetic about physical aggregates. This explains why arithmetic applies to
ordinary objects.

\subsection*{Whiteboard}

\[
\text{Number} = \text{number of things}
\]

\[
\text{``two''} = \text{all pairs}
\]

\[
2+1=3
\quad
\text{from collecting and separating objects}
\]

\section*{17. Frege's Criticism: What Are the Units?}

Frege objects that Mill's account cannot explain how units are determined.

A physical aggregate can be counted in different ways.

\[
\text{one pack of cards}
\]

\[
\text{four suits}
\]

\[
\text{fifty-two cards}
\]

The physical object alone does not determine the number. We need a concept under which
we are counting.

Likewise, a bundle of straw can be counted as:

\[
\text{one bundle}
\]

\[
\text{many straws}
\]

\[
\text{many half-straws}
\]

\[
\text{many molecules}
\]

\subsection*{Teaching Point}

Counting is not just perceiving a physical aggregate. It requires identifying units under a
concept.

\subsection*{Whiteboard}

\[
\text{Same physical stuff}
\quad \rightarrow \quad
\begin{cases}
1 \text{ pack} \\
4 \text{ suits} \\
52 \text{ cards}
\end{cases}
\]

\[
\text{Question:}
\quad
\text{What determines the units?}
\]

\section*{18. Frege's Criticism: Large Numbers}

Frege also presses Mill on large numbers.

What experience confirms a sum like:

\[
1{,}234{,}457{,}890 + 6{,}792 = 1{,}234{,}464{,}682?
\]

We do not verify this by collecting billions of objects. In fact, empirical counting would be
a terrible way to confirm such a sum. Human error would be overwhelming.

We know the sum by calculation, not by direct observation.

\subsection*{Teaching Point}

Mill's inductive account works best for small arithmetic. It struggles with large numbers,
infinity, and higher mathematics.

\subsection*{Whiteboard}

\[
\text{Small sums:}
\quad
2+1=3
\]

\[
\text{Large sums:}
\quad
1{,}234{,}457{,}890 + 6{,}792 = ?
\]

\[
\text{Do we confirm this by observation?}
\]

\section*{19. Mill on Apparent Necessity}

Mill does not deny that mathematics seems necessary. He explains this appearance
psychologically.

Basic arithmetic and geometry are confirmed by early and constant experience. We have
encountered these regularities since childhood, and we cannot imagine them failing.

\[
\text{apparent necessity}
=
\text{early and constant experience}
\]

Mathematical necessity is not genuine a priori necessity. It is the result of deeply
entrenched empirical expectation.

\subsection*{Teaching Point}

Mill replaces metaphysical necessity with psychological irresistibility formed by repeated
experience.

\subsection*{Whiteboard}

\[
\text{Why does math seem necessary?}
\]

\[
\text{Because experience has always confirmed it.}
\]

\[
\text{Necessity}
\quad \rightarrow \quad
\text{habit + imagination}
\]

\section*{20. Mill's Limits}

Mill's account has several major limits.

\[
\begin{array}{c|p{0.62\textwidth}}
\textbf{Problem} & \textbf{Why it is difficult for Mill} \\
\hline
\text{Exact geometry} &
experience gives approximations, not exact lines and points \\

\text{Large arithmetic} &
we do not confirm large sums by observation \\

\text{Mathematical induction} &
enumerative induction does not justify mathematical induction \\

\text{Infinity} &
we do not experience infinite collections \\

\text{Higher mathematics} &
set theory, analysis, and abstract mathematics exceed simple empirical construction
\end{array}
\]

\subsection*{Teaching Point}

Mill's empiricism is bold and naturalistic, but it is too weak to account for the full range
of mathematics.

\subsection*{Whiteboard}

\[
\textbf{Mill's problem:}
\]

\[
\text{experience is finite, approximate, and fallible}
\]

\[
\text{mathematics is infinite, exact, and certain?}
\]

\section*{21. Kitcher: A Modern Millian}

The chapter ends by connecting Mill to Philip Kitcher.

Kitcher keeps the empiricist spirit but improves Mill's account through idealization.

Instead of actual human construction, Kitcher appeals to ideal constructors.

\[
\text{actual human constructors}
\quad \longrightarrow \quad
\text{ideal constructors}
\]

Ideal constructors are fictional agents without ordinary human limits. They can draw
breadthless lines, collect very large aggregates, and perform idealized operations.

This allows a more sophisticated empiricist account of higher mathematics.

\subsection*{Teaching Point}

Kitcher shows how Mill's basic project can be modernized: mathematics is still connected
to human constructive activity, but through idealized activity rather than actual observation.

\subsection*{Whiteboard}

\[
\textbf{Mill:}
\quad
\text{actual experience}
\]

\[
\textbf{Kitcher:}
\quad
\text{idealized construction}
\]

\[
\text{empiricism + idealization}
\]

\section*{Conclusion: Kant and Mill as Near Opposites}

Kant and Mill agree that mathematics is not merely verbal or analytic.

Both treat mathematics as substantive.

But they disagree radically about its epistemology.

\[
\begin{array}{c|c}
\textbf{Kant} & \textbf{Mill} \\
\hline
\text{synthetic a priori} & \text{synthetic a posteriori} \\
\text{pure intuition} & \text{enumerative induction} \\
\text{forms of possible experience} & \text{generalizations from experience} \\
\text{necessity preserved} & \text{necessity explained away} \\
\text{applicability through structure of experience} &
\text{applicability through empirical origin}
\end{array}
\]

\section*{Final Framing}

To teach this chapter well, present Kant and Mill as two attempts to answer the same
problem:

\[
\boxed{
\text{How can mathematics be both exact and applicable?}
}
\]

Kant protects exactness, necessity, and a priority by grounding mathematics in pure
intuition and the forms of experience. But later mathematics and physics make this view
hard to sustain.

Mill protects naturalism and applicability by making mathematics empirical and inductive.
But exactness, infinity, large numbers, and higher mathematics become difficult to explain.

\[
\boxed{
\text{Kant makes mathematics too a priori for modern science.}
}
\]

\[
\boxed{
\text{Mill makes mathematics too empirical for mathematical exactness.}
}
\]

\section*{Discussion Questions}

\begin{enumerate}[leftmargin=2em]
    \item Why does mathematics create a problem for empiricism?
    \item Why does mathematics create a problem for rationalism?
    \item Is \(7+5=12\) analytic or synthetic?
    \item What does Kant mean by saying that mathematics constructs concepts?
    \item How does pure intuition explain the applicability of mathematics?
    \item Does non-Euclidean geometry refute Kant's philosophy of mathematics?
    \item Is Mill right that mathematics is based on experience?
    \item Can arithmetic be understood as generalization from collecting and separating objects?
    \item What determines the units when we count?
    \item Does mathematics seem necessary only because of habit and repeated experience?
    \item Is Kitcher's appeal to ideal constructors a genuine improvement on Mill?
\end{enumerate}

\section*{Key Terms}

\begin{description}[leftmargin=2em]
    \item[Rationalism] The view that reason is a fundamental source of knowledge.

    \item[Empiricism] The view that sense experience is the primary source of knowledge.

    \item[Analytic judgement] For Kant, a judgement in which the predicate concept is contained in the subject concept.

    \item[Synthetic judgement] For Kant, a judgement in which the predicate concept adds something not already contained in the subject concept.

    \item[A priori] Knowable independently of particular experience.

    \item[A posteriori] Knowable through experience.

    \item[Synthetic a priori] A judgement that extends knowledge but is knowable independently of particular experience.

    \item[Intuition] For Kant, a singular representation of an object; in mathematics, tied to construction in space and time.

    \item[Pure intuition] A priori intuition of the forms of possible experience, especially space and time.

    \item[Construction of concepts] Kant's account of mathematical cognition, in which a concept is exhibited in intuition.

    \item[Transcendentalism] The view that philosophy investigates the conditions or limits of possible experience prior to empirical science.

    \item[Naturalism] The view that philosophy is continuous with science and does not stand above it as first philosophy.

    \item[Verbal proposition] Mill's term for a proposition true by definition and lacking genuine factual content.

    \item[Real proposition] Mill's term for a proposition with factual content about the world.

    \item[Enumerative induction] Inference from observed cases to a generalization.

    \item[Limit concept] A concept approached by idealizing a sequence of empirical approximations, such as thinner and straighter lines.

    \item[Ideal constructor] Kitcher's fictional idealized agent capable of mathematical constructions beyond actual human limitations.
\end{description}

\end{document}